\subsection{Finite scan identities}\label{sec:word-calculus}
It is useful to distinguish the comparison of finite words from the
insertion of support. For words $R,W$ over any linearly ordered alphabet,
let $\res(R,W)$ be the root suffix left by the three comparison rules:
delete the smaller head, or both equal heads, and stop when one word is
empty. A missing root leaves no residual parent. All parent chains below
are stopped at that point or at a parentless head.

\begin{lemma}[Absorption and residuals]\label{lem:word}
Scanning successive target words composes:
\[
 \res(R,UV)=\res(\res(R,U),V).
\]
If $x<y$, the target word $xy$ has the same action as $y$ on every
root word. Consequently every target word has a unique weakly decreasing
normal form obtained by repeatedly removing a letter that precedes a
strictly larger letter. The residual action depends only on this normal
form. A lower same-parent root inserted immediately before a higher root
does not change the parent of the first residual head.
\end{lemma}
\begin{proof}
The first identity is the scan with its target pointer paused between
$U$ and $V$. For $xy$, roots below $x$ are removed in either scan. If
the next root equals $x$, consuming it before scanning $y$ has the same
effect as removing it as a smaller root during the $y$ scan. If the next
root lies strictly between $x$ and $y$, the first target is removed and
that root is then removed by $y$. If it is at least $y$, the first
target disappears without affecting it. These are all cases.

A stack algorithm gives the normal form: read from left to right,
remove terminal stack letters smaller than the new letter, and append
the new letter. The stack is weakly decreasing; the preceding absorption
identity proves that its residual action is unchanged. Its letters are
exactly the right-to-left weak maxima, with their multiplicities, so it
is independent of the order of removals.

For the last assertion compare the scan with and without the inserted
root $V<R$. If the two residual suffixes differ, one starts at $V$ and
the other at $R$. Their first parents agree. If $V$ is removed, the
pointers agree again; if a target removes $R$ by equality, it has already
removed $V$. This proves parent neutrality, not equality of the two
residual words. Repeating it handles a finite virtual group before one
right carrier.
\end{proof}

In particular, with $0$ the empty height, if scanning $W$ against
$U0$ leaves exactly its final $0$, then
\begin{equation}\label{eq:2-1}
             \nf(U)\le\nf(W)<\nf(U)0,
             \qquad \nf(W)=\nf(U).                 
\end{equation}
The equality follows because no weakly decreasing finite word lies
strictly between a word and that word followed by the least letter.
This identity handles a bare edge newly created by Cut. It does not
require the letters of $U$ to be bare.

For completed scans, record also a Boolean tag. A strict targetOnly
event begins a coverage block only at the least legal candidate; further
targets may be consumed by that same root while the tag persists. A
rootOnly or nonfinal both event resets the tag. This is a statement about
actual positions in the completed word, so repeated equal heights are
not identified as one occurrence.

\begin{lemma}[Composition of qualified scans]\label{lem:composition}
Suppose the legal root classes are nested, and their least-candidate
operators satisfy $C_aC_b=C_a$. For three compatible completed prefixes,
the scan from the first to the third is the composition of their actual
consumer maps through the second. If the two endpoint heights are equal,
every segment of the corresponding consumer chain is a both event.
\end{lemma}
\begin{proof}
Partition the middle word into consecutive blocks with one consumer in
the first word, and the last word into consecutive blocks with one
consumer in the middle word. Follow the first unprocessed target in
the last word. Let $B$ be its middle consumer and $A$ the first-word
consumer of $B$. Every first-word root before $A$ has either been removed
by an earlier middle target or is strictly below the first target at
which its block could start. It is likewise removed or passed in the
composed scan. If $A$ opens coverage, minimality gives
$A=C_a(C_b(T))=C_a(T)$ at the opening target $T$. Thus no earlier legal
root can replace $A$.

While the middle consumer remains $B$, the same $A$ continues to cover
the corresponding last-word block. When $B$ changes, compare the new
middle head with $A$: a smaller head remains in the block, equality
ends it by both, and a larger head removes $A$ before the next block.
These are exactly the three direct-scan events. This is a finite
induction on the two moving pointers, and does not require either word
to be sorted. At every consumption the root height is at least the
target height. If endpoint heights agree, every intermediate inequality
is an equality. The actual events are therefore both, as asserted.
\end{proof}

The equality is about completed compatible prefixes. Removing a root
that is still needed by a retained target destroys its premise.
Conversely, after deleting an unused virtual package, replay and minimal
support in Lemma~\ref{lem:replay} re-establish the same premise. This is
why support pruning is performed through birth records and actual
consumers, rather than by erasing all virtual entries at once.
